% =============================================================================
% BAB II TINJAUAN PUSTAKA -- teks panduan templat Sain-Tek-Kes 2026
% =============================================================================

\chapter{TINJAUAN PUSTAKA (OPSIONAL)}

Tinjauan pustaka berisi telaah/ulasan atas pustaka-pustaka yang relevan dengan
topik tugas akhir untuk mendapatkan informasi yang lengkap terkait kemajuan
ipteks yang telah diketahui sampai yang terkini (\textit{state of the art}).
Hal ini untuk meyakinkan pembaca bahwa tugas akhir yang dilaporkan adalah
pengetahuan baru yang lebih maju dari pengetahuan sebelumnya. Pustaka yang
digunakan dalam bab ini ialah acuan primer, diutamakan artikel jurnal dan paten
yang relevan dengan bidang yang diteliti, terkini, dan asli. Pustaka acuan harus
kredibel, dan mutakhir (setidaknya 80\% dalam 10 tahun terakhir). Diktat dan
buku ajar tidak termasuk acuan primer. Referensi lama berupa teori umum/mapan
tetap dapat digunakan. Tinjauan pustaka ditulis dengan ketentuan, jumlah halaman
bab tersebut tidak melebihi 10\% dari total halaman bagian utama naskah dan
tidak melebihi jumlah halaman Hasil dan Pembahasan. Tinjauan pustaka memuat
telaah singkat, jelas, dan sistematis tentang kerangka teoretis, kerangka pikir,
temuan, postulat-postulat, prinsip, asumsi, dan hasil penelitian yang relevan
yang melandasi topik tugas akhir atau gagasan guna menggali pemahaman mengenai
masalah kajian tugas akhir dan pemecahan masalahnya. Dalam penelitian bidang
ilmu sosial dan ekonomi, tinjauan pustaka menjadi dasar penyusunan kerangka
analisis baru dan hipotesis baru dalam topik karya ilmiah tersebut.

\section[Judul Subbab]{Judul Subbab (Kata dalam judul diawali huruf kapital dan dicetak tebal)}

Uraian dengan deskripsi untuk judul subbab 1 \dotfill

\noindent\dotfill

\noindent\dotfill

\begin{enumerate}[label=\alph*,leftmargin=0.75cm,labelwidth=0.5cm,
  labelsep=0.25cm,align=left,itemsep=0pt,parsep=0pt,topsep=0pt,partopsep=0pt]
  \item \ldots
  \item \ldots
  \item \ldots
  \begin{enumerate}[label=\arabic*),leftmargin=0.75cm,labelwidth=0.5cm,
    labelsep=0.25cm,align=left,itemsep=0pt,parsep=0pt,topsep=0pt,partopsep=0pt]
    \item \ldots
    \item \ldots
    \item \ldots
  \end{enumerate}
\end{enumerate}

\subsection{Judul Sub-subbab (Kata dalam judul diawali huruf kapital dan dicetak tidak tebal)}

\begin{ipbsubsubcontent}
Berikut adalah contoh uraian dengan deskripsi pada sub-subbab. Pada sub-subbab
ini posisi paragraf lebih menjorok 0,5 cm dari paragraf di subbab. \dotfill

\noindent\dotfill

\noindent\dotfill

\noindent\dotfill

\begin{enumerate}[label=\alph*),leftmargin=0.75cm,labelwidth=0.5cm,
  labelsep=0.25cm,align=left,itemsep=0pt,parsep=0pt,topsep=0pt,partopsep=0pt]
  \item \ldots
  \item \ldots
  \item \ldots
  \begin{enumerate}[label=(\arabic*),leftmargin=0.75cm,labelwidth=0.5cm,
    labelsep=0.25cm,align=left,itemsep=0pt,parsep=0pt,topsep=0pt,partopsep=0pt]
    \item \ldots
    \item \ldots
    \item \ldots
  \end{enumerate}
\end{enumerate}
\end{ipbsubsubcontent}

\subsection{Judul Sub-subbab (Kata dalam judul diawali huruf kapital dan dicetak tidak tebal)}

\begin{ipbsubsubcontent}
Berikut adalah contoh uraian dengan deskripsi pada sub-subbab. Pada sub-subbab
ini posisi paragraf lebih menjorok 0,5 cm dari paragraf di subbab. \dotfill

\noindent\dotfill
\end{ipbsubsubcontent}

\section[Judul Subbab]{Judul Subbab (Kata dalam judul diawali huruf kapital dan dicetak tebal)}

Berikut adalah contoh uraian dengan deskripsi pada subbab. \dotfill

\noindent\dotfill

\noindent\dotfill

\section[Judul Subbab]{Judul Subbab (Kata dalam judul diawali huruf kapital dan dicetak tebal)}

Berikut adalah contoh uraian dengan deskripsi pada subbab. \dotfill

\noindent\dotfill

\noindent\dotfill

\noindent\dotfill

\noindent\dotfill\ (Gambar~\ref{fig:mukosa-trenggiling}).

\begin{figure}[htbp]
  \centering
  \vspace{0.25cm}
  \fbox{\parbox[c][9.8cm][c]{0.92\textwidth}{\centering
    Sisipkan mikrograf atau gambar hasil penelitian di sini}}
  \caption[Mikrograf mukosa lambung trenggiling
    (\textit{Manis javanica})]{Mikrograf mukosa lambung trenggiling
    (\textit{Manis javanica}). (A) Seluruh permukaan mukosa lambung
    trenggiling dilapisi oleh epitel pipih banyak-lapis yang mengalami
    keratinisasi. (B) Pada lubang kelenjar, epitel berubah menjadi epitel
    silindris (tanda panah). Lapisan epitel mukosa terdiri atas: korneum (a),
    granulosum (tanda panah), spinosum (b), dan germinativum (c). Lapis
    submucosa (d) terletak di profundal lapis germinativum. Pewarnaan HE,
    skala = 50~$\upmu$m. (dimodifikasi dari Nisa' \textit{et al.} (2010),
    \textit{Anatomia Histologia Embryologia}. 39:432--439).}
  \label{fig:mukosa-trenggiling}
\end{figure}
